\documentclass[12pt]{article}
\usepackage[UTF8]{inputenc}
\usepackage{xeCJK}
\setCJKmainfont{Sarabun}
\usepackage{enumitem}
\usepackage{geometry}
\geometry{a4paper, margin=2.5cm}
\usepackage{setspace}
\onehalfspacing

\begin{document}

\begin{center}
{\Large\textbf{กลศาสตร์ควอนตัมเบื้องต้น}}\\[6pt]
{\normalsize วิชาฟิสิกส์ ม.ปลาย บทที่ 26}\\[4pt]
{\footnotesize ทั้งหมด 25 ข้อ}
\end{center}
\hrule
\bigskip

\section*{คำถาม in class}
\begin{enumerate}[label=\arabic*., itemsep=0.5\baselineskip, topsep=0.5\baselineskip, leftmargin=*]
\item de Broglie เสนอว่าอนุภาคทุกชนิดมีสมบัติเป็นได้ทั้งอนุภาคและคลื่น เรียกคลื่นนี้ว่าอะไร?
\begin{enumerate}[label=\alph*)., itemsep=0.3\baselineskip, topsep=0.3\baselineskip]
\item ก\ 
\begin{minipage}[t]{0.9\linewidth}
\raggedright คลื่นแม่เหล็กไฟฟ้า
\end{minipage}
\item ข\ 
\begin{minipage}[t]{0.9\linewidth}
\raggedright คลื่นอิเล็กตรออน
\end{minipage}
\item ค\ 
\begin{minipage}[t]{0.9\linewidth}
\raggedright คลื่นของ de Broglie หรือคลื่นนำความน่าจะเป็น
\end{minipage}
\item ง\ 
\begin{minipage}[t]{0.9\linewidth}
\raggedright คลื่นเสียง
\end{minipage}
\end{enumerate}
\vspace{-0.3\baselineskip}
\begin{small}
\textit{คำตอบ: ค}
\end{small}

\item ความยาวคลื่น de Broglie ของอิเล็กตรออนที่เคลื่อนที่ด้วยความเร็ว v มีสูตรอย่างไร?
\begin{enumerate}[label=\alph*)., itemsep=0.3\baselineskip, topsep=0.3\baselineskip]
\item ก\ 
\begin{minipage}[t]{0.9\linewidth}
\raggedright λ = hv
\end{minipage}
\item ข\ 
\begin{minipage}[t]{0.9\linewidth}
\raggedright λ = h/(mv)
\end{minipage}
\item ค\ 
\begin{minipage}[t]{0.9\linewidth}
\raggedright λ = mv/h
\end{minipage}
\item ง\ 
\begin{minipage}[t]{0.9\linewidth}
\raggedright λ = hf
\end{minipage}
\end{enumerate}
\vspace{-0.3\baselineskip}
\begin{small}
\textit{คำตอบ: ข}
\end{small}

\item ปรากฏการณ์โฟโตอิเล็กทริก (Photoelectric effect) เกิดขึ้นเมื่อใด?
\begin{enumerate}[label=\alph*)., itemsep=0.3\baselineskip, topsep=0.3\baselineskip]
\item ก\ 
\begin{minipage}[t]{0.9\linewidth}
\raggedright เมื่อแสงตกกระทบโลหะและอิเล็กตรออนหลุดออกมา
\end{minipage}
\item ข\ 
\begin{minipage}[t]{0.9\linewidth}
\raggedright เมื่ออิเล็กตรออนชนอะตอม
\end{minipage}
\item ค\ 
\begin{minipage}[t]{0.9\linewidth}
\raggedright เมื่อโลหะร้อนมาก
\end{minipage}
\item ง\ 
\begin{minipage}[t]{0.9\linewidth}
\raggedright เมื่อแสงสะท้อนจากกระจก
\end{minipage}
\end{enumerate}
\vspace{-0.3\baselineskip}
\begin{small}
\textit{คำตอบ: ก}
\end{small}

\item สมการโฟโตอิเล็กทริกของไอน์สไตน์คือข้อใด?
\begin{enumerate}[label=\alph*)., itemsep=0.3\baselineskip, topsep=0.3\baselineskip]
\item ก\ 
\begin{minipage}[t]{0.9\linewidth}
\raggedright E = mc²
\end{minipage}
\item ข\ 
\begin{minipage}[t]{0.9\linewidth}
\raggedright hf = φ + K\_max
\end{minipage}
\item ค\ 
\begin{minipage}[t]{0.9\linewidth}
\raggedright E = pc
\end{minipage}
\item ง\ 
\begin{minipage}[t]{0.9\linewidth}
\raggedright hf = mc²
\end{minipage}
\end{enumerate}
\vspace{-0.3\baselineskip}
\begin{small}
\textit{คำตอบ: ข}
\end{small}

\item หลักความไม่แน่นอนของไฮเซนเบิร์ก (Heisenberg Uncertainty Principle) บอกว่าอะไรไม่สามารถวัดได้พร้อมกันอย่างแม่นยำ?
\begin{enumerate}[label=\alph*)., itemsep=0.3\baselineskip, topsep=0.3\baselineskip]
\item ก\ 
\begin{minipage}[t]{0.9\linewidth}
\raggedright ตำแหน่งและความเร็วของอิเล็กตรออน
\end{minipage}
\item ข\ 
\begin{minipage}[t]{0.9\linewidth}
\raggedright มวลและประจุของอิเล็กตรออน
\end{minipage}
\item ค\ 
\begin{minipage}[t]{0.9\linewidth}
\raggedright อุณหภูมิและความดัน
\end{minipage}
\item ง\ 
\begin{minipage}[t]{0.9\linewidth}
\raggedright พลังงานและเวลา
\end{minipage}
\end{enumerate}
\vspace{-0.3\baselineskip}
\begin{small}
\textit{คำตอบ: ก}
\end{small}

\item ฟังก์ชันคลื่น Ψ (psi) ในสมการ Schrödinger บอกอะไร?
\begin{enumerate}[label=\alph*)., itemsep=0.3\baselineskip, topsep=0.3\baselineskip]
\item ก\ 
\begin{minipage}[t]{0.9\linewidth}
\raggedright ตำแหน่งที่แน่นอนของอิเล็กตรออน
\end{minipage}
\item ข\ 
\begin{minipage}[t]{0.9\linewidth}
\raggedright ความน่าจะเป็นในการพบอิเล็กตรออน
\end{minipage}
\item ค\ 
\begin{minipage}[t]{0.9\linewidth}
\raggedright พลังงานของอิเล็กตรออน
\end{minipage}
\item ง\ 
\begin{minipage}[t]{0.9\linewidth}
\raggedright ความเร็วของอิเล็กตรออน
\end{minipage}
\end{enumerate}
\vspace{-0.3\baselineskip}
\begin{small}
\textit{คำตอบ: ข}
\end{small}

\item ค่าคงที่พลังค์ h มีค่าเท่าใด?
\begin{enumerate}[label=\alph*)., itemsep=0.3\baselineskip, topsep=0.3\baselineskip]
\item ก\ 
\begin{minipage}[t]{0.9\linewidth}
\raggedright 6.626 × 10⁻³⁴ J·s
\end{minipage}
\item ข\ 
\begin{minipage}[t]{0.9\linewidth}
\raggedright 6.626 × 10⁻³⁴ J/s
\end{minipage}
\item ค\ 
\begin{minipage}[t]{0.9\linewidth}
\raggedright 1.055 × 10⁻³⁴ J·s
\end{minipage}
\item ง\ 
\begin{minipage}[t]{0.9\linewidth}
\raggedright 3.14 × 10⁻³⁴ J·s
\end{minipage}
\end{enumerate}
\vspace{-0.3\baselineskip}
\begin{small}
\textit{คำตอบ: ก}
\end{small}

\item ทำไมอนุภาคมวลใหญ่ในชีวิตประจำวันจึงไม่แสดงสมบัติคลื่น?
\begin{enumerate}[label=\alph*)., itemsep=0.3\baselineskip, topsep=0.3\baselineskip]
\item ก\ 
\begin{minipage}[t]{0.9\linewidth}
\raggedright เพราะมวลเป็นศูนย์
\end{minipage}
\item ข\ 
\begin{minipage}[t]{0.9\linewidth}
\raggedright เพราะความยาวคลื่น de Broglie น้อยมากจนวัดไม่ได้
\end{minipage}
\item ค\ 
\begin{minipage}[t]{0.9\linewidth}
\raggedright เพราะอนุภาคมวลใหญ่เคลื่อนที่ช้ามาก
\end{minipage}
\item ง\ 
\begin{minipage}[t]{0.9\linewidth}
\raggedright เพราะไม่มีประจุ
\end{minipage}
\end{enumerate}
\vspace{-0.3\baselineskip}
\begin{small}
\textit{คำตอบ: ข}
\end{small}

\item ข้อใดไม่ใช่ปัจจัยที่มีผลต่อพลังงานจลน์สูงสุดของอิเล็กตรออนในปรากฏการณ์โฟโตอิเล็กทริก?
\begin{enumerate}[label=\alph*)., itemsep=0.3\baselineskip, topsep=0.3\baselineskip]
\item ก\ 
\begin{minipage}[t]{0.9\linewidth}
\raggedright ความถี่ของแสง
\end{minipage}
\item ข\ 
\begin{minipage}[t]{0.9\linewidth}
\raggedright ความเข้มของแสง
\end{minipage}
\item ค\ 
\begin{minipage}[t]{0.9\linewidth}
\raggedright ค่างานของโลหะ
\end{minipage}
\item ง\ 
\begin{minipage}[t]{0.9\linewidth}
\raggedright ชนิดของโลหะ
\end{minipage}
\end{enumerate}
\vspace{-0.3\baselineskip}
\begin{small}
\textit{คำตอบ: ข}
\end{small}

\item ไม่สามารถวัดตำแหน่งและโมเมนตัมของอิเล็กตรออนได้อย่างแม่นยำพร้อมกัน เพราะอะไร?
\begin{enumerate}[label=\alph*)., itemsep=0.3\baselineskip, topsep=0.3\baselineskip]
\item ก\ 
\begin{minipage}[t]{0.9\linewidth}
\raggedright เครื่องมือวัดไม่ละเอียดพอ
\end{minipage}
\item ข\ 
\begin{minipage}[t]{0.9\linewidth}
\raggedright การวัดทำให้ระบบเปลี่ยนแปลง
\end{minipage}
\item ค\ 
\begin{minipage}[t]{0.9\linewidth}
\raggedright เพราะอิเล็กตรออนมีสมบัติคลื่น
\end{minipage}
\item ง\ 
\begin{minipage}[t]{0.9\linewidth}
\raggedright เพราะอิเล็กตรออนเล็กเกินไป
\end{minipage}
\end{enumerate}
\vspace{-0.3\baselineskip}
\begin{small}
\textit{คำตอบ: ค}
\end{small}

\item |Ψ|² dx ในกลศาสตร์ควอนตัมหมายถึงอะไร?
\begin{enumerate}[label=\alph*)., itemsep=0.3\baselineskip, topsep=0.3\baselineskip]
\item ก\ 
\begin{minipage}[t]{0.9\linewidth}
\raggedright พลังงานของคลื่น
\end{minipage}
\item ข\ 
\begin{minipage}[t]{0.9\linewidth}
\raggedright ความน่าจะเป็นในการพบอนุภาคในช่วง dx
\end{minipage}
\item ค\ 
\begin{minipage}[t]{0.9\linewidth}
\raggedright ความยาวคลื่น
\end{minipage}
\item ง\ 
\begin{minipage}[t]{0.9\linewidth}
\raggedright โมเมนตัมของอนุภาค
\end{minipage}
\end{enumerate}
\vspace{-0.3\baselineskip}
\begin{small}
\textit{คำตอบ: ข}
\end{small}

\item อิเล็กตรออนเคลื่อนที่ด้วยความเร็ว 1.0 × 10⁶ m/s จะมีความยาวคลื่น de Broglie เท่าใด? (h = 6.63 × 10⁻³⁴ J·s, mₑ = 9.11 × 10⁻³¹ kg)
\begin{enumerate}[label=\alph*)., itemsep=0.3\baselineskip, topsep=0.3\baselineskip]
\item ก\ 
\begin{minipage}[t]{0.9\linewidth}
\raggedright 7.3 × 10⁻¹⁰ m
\end{minipage}
\item ข\ 
\begin{minipage}[t]{0.9\linewidth}
\raggedright 7.3 × 10⁻⁹ m
\end{minipage}
\item ค\ 
\begin{minipage}[t]{0.9\linewidth}
\raggedright 7.3 × 10⁻⁸ m
\end{minipage}
\item ง\ 
\begin{minipage}[t]{0.9\linewidth}
\raggedright 7.3 × 10⁻⁷ m
\end{minipage}
\end{enumerate}
\vspace{-0.3\baselineskip}
\begin{small}
\textit{คำตอบ: ข}
\end{small}

\item แสงความถี่ f ตกกระทบโลหะที่มีค่างาน φ ถ้า hf > φ จะเกิดอะไรขึ้น?
\begin{enumerate}[label=\alph*)., itemsep=0.3\baselineskip, topsep=0.3\baselineskip]
\item ก\ 
\begin{minipage}[t]{0.9\linewidth}
\raggedright อิเล็กตรออนหลุดออกมาพร้อมพลังงานจลน์ K = hf - φ
\end{minipage}
\item ข\ 
\begin{minipage}[t]{0.9\linewidth}
\raggedright อิเล็กตรออนหลุดออกมาพร้อมพลังงานจลน์ K = hf
\end{minipage}
\item ค\ 
\begin{minipage}[t]{0.9\linewidth}
\raggedright อิเล็กตรออนไม่หลุดออกมา
\end{minipage}
\item ง\ 
\begin{minipage}[t]{0.9\linewidth}
\raggedright โลหะเปล่งแสง
\end{minipage}
\end{enumerate}
\vspace{-0.3\baselineskip}
\begin{small}
\textit{คำตอบ: ก}
\end{small}

\item แสงสีแดง (λ = 650 nm) กับแสงสีม่วง (λ = 400 nm) อันใดมีพลังงานโฟตอนมากกว่า?
\begin{enumerate}[label=\alph*)., itemsep=0.3\baselineskip, topsep=0.3\baselineskip]
\item ก\ 
\begin{minipage}[t]{0.9\linewidth}
\raggedright แสงสีแดง
\end{minipage}
\item ข\ 
\begin{minipage}[t]{0.9\linewidth}
\raggedright แสงสีม่วง
\end{minipage}
\item ค\ 
\begin{minipage}[t]{0.9\linewidth}
\raggedright เท่ากัน
\end{minipage}
\item ง\ 
\begin{minipage}[t]{0.9\linewidth}
\raggedright ขึ้นอยู่กับความเข้ม
\end{minipage}
\end{enumerate}
\vspace{-0.3\baselineskip}
\begin{small}
\textit{คำตอบ: ข}
\end{small}

\item ถ้า Δx มีค่าน้อยมาก ความไม่แน่นอนของโมเมนตัม Δp จะเป็นอย่างไร?
\begin{enumerate}[label=\alph*)., itemsep=0.3\baselineskip, topsep=0.3\baselineskip]
\item ก\ 
\begin{minipage}[t]{0.9\linewidth}
\raggedright Δp มีค่าน้อย
\end{minipage}
\item ข\ 
\begin{minipage}[t]{0.9\linewidth}
\raggedright Δp มีค่ามาก
\end{minipage}
\item ค\ 
\begin{minipage}[t]{0.9\linewidth}
\raggedright Δp = 0
\end{minipage}
\item ง\ 
\begin{minipage}[t]{0.9\linewidth}
\raggedright Δp คงที่
\end{minipage}
\end{enumerate}
\vspace{-0.3\baselineskip}
\begin{small}
\textit{คำตอบ: ข}
\end{small}

\item ในกลศาสตร์ควอนตัม อิเล็กตรออนไม่ได้เคลื่อนที่เป็นวงโคจรเหมือนดาวเคราะห์รอบดวงอาทิตย์ แต่มีสมบัติอย่างไร?
\begin{enumerate}[label=\alph*)., itemsep=0.3\baselineskip, topsep=0.3\baselineskip]
\item ก\ 
\begin{minipage}[t]{0.9\linewidth}
\raggedright อิเล็กตรออนหยุดนิ่ง
\end{minipage}
\item ข\ 
\begin{minipage}[t]{0.9\linewidth}
\raggedright อิเล็กตรออนมีโอกาสพบได้ในบริเวณที่เรียกว่า atomic orbital
\end{minipage}
\item ค\ 
\begin{minipage}[t]{0.9\linewidth}
\raggedright อิเล็กตรออนเคลื่อนที่เร็วมากจนวัดไม่ได้
\end{minipage}
\item ง\ 
\begin{minipage}[t]{0.9\linewidth}
\raggedright อิเล็กตรออนอยู่ในสถานะของแข็ง
\end{minipage}
\end{enumerate}
\vspace{-0.3\baselineskip}
\begin{small}
\textit{คำตอบ: ข}
\end{small}

\item การทดลองใดเป็นหลักฐานยืนยันสมมติฐานของ de Broglie?
\begin{enumerate}[label=\alph*)., itemsep=0.3\baselineskip, topsep=0.3\baselineskip]
\item ก\ 
\begin{minipage}[t]{0.9\linewidth}
\raggedright การทดลองเรืองแสง
\end{minipage}
\item ข\ 
\begin{minipage}[t]{0.9\linewidth}
\raggedright การเลี้ยวเบนของอิเล็กตรออนโดยผลึก
\end{minipage}
\item ค\ 
\begin{minipage}[t]{0.9\linewidth}
\raggedright การกระเจิงของอนุภาคแอลฟา
\end{minipage}
\item ง\ 
\begin{minipage}[t]{0.9\linewidth}
\raggedright การแผ่รังสีของวัตถุดำ
\end{minipage}
\end{enumerate}
\vspace{-0.3\baselineskip}
\begin{small}
\textit{คำตอบ: ข}
\end{small}

\item ถ้าเพิ่มความถี่ของแสงที่ตกกระทบโลหะในขณะที่ความเข้มคงที่ ผลที่เกิดขึ้นคืออะไร?
\begin{enumerate}[label=\alph*)., itemsep=0.3\baselineskip, topsep=0.3\baselineskip]
\item ก\ 
\begin{minipage}[t]{0.9\linewidth}
\raggedright จำนวนอิเล็กตรออนที่หลุดลดลง แต่พลังงานจลน์เพิ่มขึ้น
\end{minipage}
\item ข\ 
\begin{minipage}[t]{0.9\linewidth}
\raggedright จำนวนอิเล็กตรออนเพิ่มขึ้นและพลังงานจลน์เพิ่มขึ้น
\end{minipage}
\item ค\ 
\begin{minipage}[t]{0.9\linewidth}
\raggedright จำนวนอิเล็กตรออนคงที่ พลังงานจลน์ลดลง
\end{minipage}
\item ง\ 
\begin{minipage}[t]{0.9\linewidth}
\raggedright ไม่มีผลอะไร
\end{minipage}
\end{enumerate}
\vspace{-0.3\baselineskip}
\begin{small}
\textit{คำตอบ: ก}
\end{small}

\item แสง UV (f = 1.5 × 10¹⁵ Hz) ตกกระทบโลหะที่มีค่างาน φ = 3.0 eV อิเล็กตรออนจะมีพลังงานจลน์สูงสุดเท่าใด?
\begin{enumerate}[label=\alph*)., itemsep=0.3\baselineskip, topsep=0.3\baselineskip]
\item ก\ 
\begin{minipage}[t]{0.9\linewidth}
\raggedright 1.2 eV
\end{minipage}
\item ข\ 
\begin{minipage}[t]{0.9\linewidth}
\raggedright 2.2 eV
\end{minipage}
\item ค\ 
\begin{minipage}[t]{0.9\linewidth}
\raggedright 3.2 eV
\end{minipage}
\item ง\ 
\begin{minipage}[t]{0.9\linewidth}
\raggedright 4.2 eV
\end{minipage}
\end{enumerate}
\vspace{-0.3\baselineskip}
\begin{small}
\textit{คำตอบ: ค}
\end{small}

\item อิเล็กตรออนถูกจำกัดให้อยู่ในกล่องกว้าง 0.1 nm ความไม่แน่นอนของโมเมนตัมมีค่าอย่างน้อยเท่าใด?
\begin{enumerate}[label=\alph*)., itemsep=0.3\baselineskip, topsep=0.3\baselineskip]
\item ก\ 
\begin{minipage}[t]{0.9\linewidth}
\raggedright 5.27 × 10⁻²⁵ kg·m/s
\end{minipage}
\item ข\ 
\begin{minipage}[t]{0.9\linewidth}
\raggedright 5.27 × 10⁻²⁴ kg·m/s
\end{minipage}
\item ค\ 
\begin{minipage}[t]{0.9\linewidth}
\raggedright 5.27 × 10⁻²³ kg·m/s
\end{minipage}
\item ง\ 
\begin{minipage}[t]{0.9\linewidth}
\raggedright 5.27 × 10⁻²² kg·m/s
\end{minipage}
\end{enumerate}
\vspace{-0.3\baselineskip}
\begin{small}
\textit{คำตอบ: ข}
\end{small}

\item โปรตอน (m = 1.67 × 10⁻²⁷ kg) เคลื่อนที่ด้วยความเร็ว 1.0 × 10⁴ m/s มีความยาวคลื่น de Broglie เท่าใด?
\begin{enumerate}[label=\alph*)., itemsep=0.3\baselineskip, topsep=0.3\baselineskip]
\item ก\ 
\begin{minipage}[t]{0.9\linewidth}
\raggedright 4.0 × 10⁻¹¹ m
\end{minipage}
\item ข\ 
\begin{minipage}[t]{0.9\linewidth}
\raggedright 4.0 × 10⁻¹⁰ m
\end{minipage}
\item ค\ 
\begin{minipage}[t]{0.9\linewidth}
\raggedright 4.0 × 10⁻⁹ m
\end{minipage}
\item ง\ 
\begin{minipage}[t]{0.9\linewidth}
\raggedright 4.0 × 10⁻⁸ m
\end{minipage}
\end{enumerate}
\vspace{-0.3\baselineskip}
\begin{small}
\textit{คำตอบ: ก}
\end{small}

\item ค่าความถี่ขีดเริ่ม (threshold frequency) ของโลหะคืออะไร?
\begin{enumerate}[label=\alph*)., itemsep=0.3\baselineskip, topsep=0.3\baselineskip]
\item ก\ 
\begin{minipage}[t]{0.9\linewidth}
\raggedright ความถี่ต่ำสุดที่ทำให้อิเล็กตรออนหลุด
\end{minipage}
\item ข\ 
\begin{minipage}[t]{0.9\linewidth}
\raggedright ความถี่สูงสุดที่ทำให้อิเล็กตรออนหลุด
\end{minipage}
\item ค\ 
\begin{minipage}[t]{0.9\linewidth}
\raggedright ความถี่ที่ทำให้อิเล็กตรออนหลุดทั้งหมด
\end{minipage}
\item ง\ 
\begin{minipage}[t]{0.9\linewidth}
\raggedright ความถี่ที่ไม่ทำให้อิเล็กตรออนหลุด
\end{minipage}
\end{enumerate}
\vspace{-0.3\baselineskip}
\begin{small}
\textit{คำตอบ: ก}
\end{small}

\item ในกลศาสตร์ควอนตัม ถ้าวัดตำแหน่งของอิเล็กตรออนหลายครั้งในสภาพเดียวกัน ผลการวัดจะเป็นอย่างไร?
\begin{enumerate}[label=\alph*)., itemsep=0.3\baselineskip, topsep=0.3\baselineskip]
\item ก\ 
\begin{minipage}[t]{0.9\linewidth}
\raggedright ได้ผลเหมือนกันทุกครั้ง
\end{minipage}
\item ข\ 
\begin{minipage}[t]{0.9\linewidth}
\raggedright ได้ผลกระจายตัว แต่มีค่าเฉลี่ยตาม |Ψ|²
\end{minipage}
\item ค\ 
\begin{minipage}[t]{0.9\linewidth}
\raggedright ไม่สามารถวัดได้
\end{minipage}
\item ง\ 
\begin{minipage}[t]{0.9\linewidth}
\raggedright ได้ผลที่แปลกประหลาดเสมอ
\end{minipage}
\end{enumerate}
\vspace{-0.3\baselineskip}
\begin{small}
\textit{คำตอบ: ข}
\end{small}

\item แสงเลเซอร์ความเข้มสูงมากทำให้เกิดปรากฏการณ์โฟโตอิเล็กทริกได้ง่ายกว่าแสงปกติ เพราะอะไร?
\begin{enumerate}[label=\alph*)., itemsep=0.3\baselineskip, topsep=0.3\baselineskip]
\item ก\ 
\begin{minipage}[t]{0.9\linewidth}
\raggedright เพราะแสงเลเซอร์มีความถี่สูงกว่าเสมอ
\end{minipage}
\item ข\ 
\begin{minipage}[t]{0.9\linewidth}
\raggedright เพราะแสงเลเซอร์มีความเข้มสูงมากต่อหน่วยพื้นที่ ทำให้มีโฟตอนจำนวนมาก
\end{minipage}
\item ค\ 
\begin{minipage}[t]{0.9\linewidth}
\raggedright เพราะแสงเลเซอร์มีโฟตอนพลังงานสูงกว่า
\end{minipage}
\item ง\ 
\begin{minipage}[t]{0.9\linewidth}
\raggedright เพราะแสงเลเซอร์เป็นแสงขาว
\end{minipage}
\end{enumerate}
\vspace{-0.3\baselineskip}
\begin{small}
\textit{คำตอบ: ข}
\end{small}

\item อะตอมมีขนาดประมาณ 0.1 nm ถ้าอิเล็กตรออนถูกกักไว้ในอะตอม ความไม่แน่นอนของความเร็วจะมีค่าอย่างน้อยเท่าใด?
\begin{enumerate}[label=\alph*)., itemsep=0.3\baselineskip, topsep=0.3\baselineskip]
\item ก\ 
\begin{minipage}[t]{0.9\linewidth}
\raggedright ≈ 10⁶ m/s
\end{minipage}
\item ข\ 
\begin{minipage}[t]{0.9\linewidth}
\raggedright ≈ 10⁵ m/s
\end{minipage}
\item ค\ 
\begin{minipage}[t]{0.9\linewidth}
\raggedright ≈ 10⁴ m/s
\end{minipage}
\item ง\ 
\begin{minipage}[t]{0.9\linewidth}
\raggedright ≈ 10³ m/s
\end{minipage}
\end{enumerate}
\vspace{-0.3\baselineskip}
\begin{small}
\textit{คำตอบ: ก}
\end{small}

\end{enumerate}
\end{document}